\section{Conclusion}
\label{sec:conclusion}

This paper asked one question: can tensor cores be applied to ALS matrix factorization? Tensor cores sit inside modern GPUs, including ordinary gaming cards, and they do dense matrix multiplication very fast, but a recommendation system deals with a sparse matrix, and a sparse matrix does not feed tensor cores well. The answer built here rests on sorting. Users are sorted so that the ones who give the most ratings come first, and items so that the most popular come first. After the sorting, the dense part of the matrix gathers in one corner, and on that corner tensor cores become usable. APR-BALS wraps this idea in a complete training system: a three-tier dispatch sends the dense work to the tensor cores in FP16 with FP32 accumulation, the sparse tail goes to a register-tiled scalar FP32 kernel, and a blocked batched Cholesky solver keeps the $K^{3}$ solve stage from eating the gain. The baseline and the accelerated path share one binary, so the comparison stays fair. On an RTX~3060 the Netflix Prize dataset trains $4.6\times$ to $15.4\times$ faster than the FP32 baseline over $K = 16$--$96$, the test RMSE moves less than $0.13\%$, the results repeat on MovieLens 20M and 10M and on a Tesla T4, and the gap to multi-threaded CPU and PyTorch/cuBLAS references on the same data is about two orders of magnitude.

The method has clear borders. The supported latent dimensions are $K = 16, 32, 48, 64, 96$ only, enforced by a compile-time guard; adding another value needs a new WMMA template instantiation plus a register-spill check. At high $K$, most of the remaining time passes inside the tiled Cholesky kernel, about $57\%$ of the accelerated iteration at $K=96$, and phase-isolation probes indicate that kernel now sits near a latency-bound structural floor. The system optimizes explicit ratings and evaluates by RMSE; implicit-feedback weighting and ranking metrics are outside the current scope. Only a single GPU is used. The sparse path uses atomic additions, so an APR run repeats only to about $10^{-6}$; bitwise-reproducible training would require a deterministic reduction. And although high-$K$ training is now cheap, $K = 16$--$32$ gives the best test accuracy on these datasets, a property of the data that hits both precision paths equally.

Several of the extensions proposed below were subsequently completed and are reported in Section~\ref{ssec:extensions}: a narrower-precision solver variant enabled by switching to weighted regularization (a $23\%$ solve-stage reduction at $K=96$, with test RMSE matching the FP32 solver exactly); the removal of the two largest non-solver phases by fusing the regularizer into the solver load and relocalizing the RMSE check (a further $18\%$ off the $K=96$ wall time); a head-to-head comparison against cuMF\_als \cite{tan2016faster} on the same frozen split across five datasets, where APR-BALS wins per-iteration speed by $3.0\times$--$6.4\times$ while agreeing with cuMF's train RMSE to within $0.33\%$; and a heavy-ball momentum extrapolation of the ALS factor sequence that, orthogonally to the per-sweep kernel work, cuts the sweep count to convergence from 20--25 to 15 and the $K=96$ wall time by a further $1.30\times$ on Netflix and up to $1.65\times$ on MovieLens, at equal or better test accuracy. The comparison also settled how far the solver can be pushed: a vendor batched factorization (cuSOLVER) and a tensor-core batched-multiply formulation were both implemented and measured slower, because at these ranks the per-entity matrices are below the size at which the tensor cores factor efficiently, so the shared-memory scalar Cholesky is retained. What remains on the solver is writing the Gram matrices directly in its tiled layout so the load phase becomes a straight vectorized copy; the already-implemented dual-stream pipeline would convert any solver slimming into wall-clock gains automatically. BF16 and TF32 operand formats run on the same tensor cores, widen the exponent range, and would likely shrink the already small high-$K$ deltas; extending the dispatch to the confidence-weighted implicit-feedback objective \cite{hu2008collaborative} would broaden applicability; and RMSE could be complemented with ranking metrics.

One last point goes beyond recommendation. Many sparse workloads hide dense per-entity substructure: ridge-regression sweeps, Gaussian-process approximations, graph feature aggregation. The same tiered-precision pattern has a fair chance of working there too, and the honest-baseline habit of this work travels even further. Any acceleration claim, in any field, becomes stronger when it is measured this way.
